\section{Discusion e implicancias}

Los resultados no respaldan una narrativa unica de espiral sostenida en toda la muestra; mas bien, sugieren comportamiento dependiente de episodio y especificacion. Para politica publica, esto implica que la calibracion de reajustes del SML deberia evaluarse junto con:
\begin{itemize}
  \item composicion sectorial y de formalidad del empleo;
  \item dinamica de precios en alimentos frente a IPC general;
  \item condiciones de bienestar del hogar que afectan vulnerabilidad inflacionaria.
\end{itemize}

La evidencia de bunching alrededor de 1 SML y los diferenciales por contexto social abren agenda para estrategias complementarias: fiscalizacion laboral, politicas de productividad y focalizacion por hogares de alta exposicion.

\section{Conclusion}

Este articulo presenta una arquitectura reproducible para estudiar la relacion salario minimo--precios en Paraguay combinando evento, placebo y contexto micro. En la especificacion principal, los EI promedio son moderados (\eiGeneralMain{} pp en IPC general y \eiAlimentosMain{} pp en IPC alimentos) y estadisticamente fragiles frente a placebo. Por tanto, la contribucion principal es metodologica y de transparencia: proveer evidencia trazable, actualizable y util para dialogo tecnico.
